% !TeX root = ../main.tex
\section{Nhu cầu gọi vốn}

\textbf{1. Xác định nhu cầu vốn đầu tư:}
Dựa trên mô hình dự phóng tài chính và phân tích dòng tiền ở Phần 4, khoản thiếu hụt dòng tiền tích lũy lớn nhất (Maximum Cash Drain) của Mutux rơi vào \textbf{Quý 3 năm đầu tiên} với mức âm tích lũy là \textbf{-231 triệu VNĐ}. Mặc dù sang Quý 4 dòng tiền thuần từ hoạt động kinh doanh bắt đầu đạt dương (+\textbf{91 triệu VNĐ}), tổng khoản thiếu hụt tài chính tích lũy toàn năm 1 vẫn ở mức khoảng \textbf{140 - 156 triệu VNĐ}.

Nhằm đảm bảo an toàn tài chính, chuẩn bị quỹ dự phòng cho rủi ro vận hành (rủi ro thanh khoản cọc, rủi ro nợ xấu tài sản) và duy trì thời gian hoạt động (Runway) tối thiểu 18 tháng, Mutux xác định nhu cầu gọi vốn cho vòng Tiền hạt giống (Pre-Seed Round) là \textbf{12.000 USD} (tương đương khoảng \textbf{300.000.000 VNĐ}).

\textbf{2. Mục tiêu chiến lược của vòng gọi vốn Pre-Seed:}
\begin{itemize}
    \item \textbf{Đảm bảo Runway hoạt động:} Cung cấp nguồn vốn an toàn kéo dài ít nhất 18--24 tháng, giúp startup vượt qua giai đoạn khởi tạo (Cold-start) và duy trì bộ máy vận hành tinh gọn.
    \item \textbf{Đạt điểm hòa vốn (Break-even Point):} Tài trợ nguồn lực để thúc đẩy tốc độ tăng trưởng giao dịch, đưa nền tảng đạt điểm hòa vốn vào \textbf{năm thứ 2 (Quý 3)} khi doanh thu định kỳ đạt mức 80--100 triệu VNĐ/tháng.
    \item \textbf{Tạo đà bứt phá quy mô:} Đạt mục tiêu 10.000 người dùng đăng ký, 7.050 đơn thuê thành công và 1,11 tỷ VNĐ tổng giá trị giao dịch (GMV) trong 12 tháng đầu, làm tiền đề mở rộng các dịch vụ B2B và thương mại hóa chuyên sâu.
\end{itemize}
